\documentclass[a4paper]{article}     
\usepackage{amsfonts}
\usepackage{amsmath}
\usepackage{amssymb}
\usepackage{graphicx}
\usepackage{cancel}

\newcommand{\asa}{\Leftrightarrow} %als slechts als
\newcommand{\adR}{\Rightarrow} %als dan Rechts
\newcommand{\adL}{\Leftarrow}
\newcommand{\Lh}{\left(} %Links haakje
\newcommand{\Rh}{\right)}
\newcommand{\Lv}{\left[} %Links vierkanthaakje
\newcommand{\Rv}{\right]}
\newcommand{\La}{\left\{} %Links acolade
\newcommand{\Ra}{\right\}}
\newcommand{\Ras}{\right.} %Rechts acolade stelsel (omdat om stelsels te maken heb je een punt nodig
\newcommand{\pnR}{\rightarrow} %pijl naar Rechts
\newcommand{\limiet}{\lim\limits_}
\newcommand{\schrap}{\cancel}
\newcommand{\schrapb}{\bcancel} %backslash
\newcommand{\schrapk}{\xcancel} %kruis
\newcommand{\vervang}{\cancelto} 

\begin{document}

\noindent \large Auteur: Yoren Collier \\
\noindent \large Studierichting: Informatica\\
\noindent \large Oefeningenreeks 2D nr. 10.9\\

\medskip

\normalsize

$f(x) = \dfrac{x^2 +x -2}{x^2 -4x +3}$\\

$\limiet{x \pnR 1} f(x) = \dfrac{1+1-2}{1-4+3} = \dfrac{0}{0}$\\

We weten dat 1 een nulpunt is voor beide $t(x)$ en $n(x)$. Met $t(x)$ de (veelterm)functie van de teller en $n(x)$ de (veelterm)functie van de noemer. Hierdoor kunnen we Horner gebruiken en de teller en noemer ontbinden zodat we verder het limiet kunnen berekenen. \\

Horner voor de teller:\\

\begin{tabular}{c|ccc}
     &  $1$  &  $1$  &  $-2$  \\
  $1$&  $|$  &  $1$  &   $2$  \\
     &  $1$  &  $2$  &   $0$  
\end{tabular}

$\adR t(x) = (x-1)(x+2)$\\

Horner voor de noemer:\\

\begin{tabular}{c|ccc}
     &  $1$  &  $-4$  &  $3$  \\
  $1$&  $|$  &  $1$   &  $-3$  \\
     &  $1$  &  $-3$  &  $0$  
\end{tabular}

$\adR n(x) = (x-1)(x-3)$\\

$\limiet{x \pnR 1} f(x) = \limiet{x \pnR 1} \dfrac{\schrap{\Lh x-1 \Rh} \Lh x+2 \Rh }{\schrap{\Lh x-1 \Rh} \Lh x-3 \Rh } = \dfrac{1+2}{1-3} = \dfrac{-3}{2}$\\

\begin{thebibliography}{999}
%\bibitem{a} Referentie
\end{thebibliography}
\end{document} 